\chapter{Python の基本}
\label{chap:python_basics}

% この章で学ぶこと
\begin{goalbox}
  \begin{itemize}
    \item 数値の計算と、変数・データ型の使い方
    \item 文字列の扱い方と、キーボードからの入力
    \item リスト・タプル・辞書を使って、複数のデータをまとめる方法
    \item if 文による条件分岐と、for 文・while 文による繰り返し
    \item 関数を作って、処理をまとめる方法
  \end{itemize}
\end{goalbox}

この章では、Python でプログラムを書くための基本的な文法を学びます。
各節の最後には練習問題を用意しました。
解答例は付録\ref{app:python-answers}（\ref{sec:python-basics-answers}節）にあります。
プログラミングは、読むだけでなく、実際に書いて動かしてみることで身につきます。
ぜひ、自分の手でプログラムを書き、いろいろと書き換えて試してみてください。

\begin{point}[この章のプログラムの実行方法]
  この章のプログラムは、\ref{sec:python-run}節で学んだ方法で実行します。
  短い例は対話モード（\cmd{>>>} から始まる例）で、長い例はスクリプトファイルに保存して \cmd{python3 ファイル名} で実行してください。
  練習用のディレクトリ（たとえば \cmd{\textasciitilde/py\_practice}）を作り、その中でファイルを作っていくとよいでしょう。
\end{point}

\section{数値と計算}
\label{sec:python-basics-number}

\subsection{四則演算}

Python では、数値と演算子を組み合わせて計算ができます。
足し算は \cmd{+}、引き算は \cmd{-}、掛け算は \cmd{*}、割り算は \cmd{/} です。

\begin{terminal}[四則演算]
>>> 3 + 5
8
>>> 6 - 10
-4
>>> 6 * 8
48
>>> 26 / 4
6.5
\end{terminal}

算数と同じように、掛け算と割り算が先に計算されます。
先に計算したい部分は \cmd{( )} で囲みます。

\begin{terminal}[計算の順番]
>>> 10 - 5 / 5
9.0
>>> (10 - 5) / 5
1.0
\end{terminal}

\subsection{四則演算以外の計算}

ほかにも、次のような演算子があります。

\begin{tablebox}[四則演算以外の演算子][tab:python-basics-operators]
  \begin{tabular}{lll}
    \toprule
    演算子 & 意味 & 例 \\
    \midrule
    \cmd{**} & 累乗 & \cmd{2 ** 5} は 32 \\
    \cmd{//} & 割り算の商（小数点以下を切り捨て） & \cmd{10 // 3} は 3 \\
    \cmd{\%} & 割り算のあまり & \cmd{10 \% 3} は 1 \\
    \bottomrule
  \end{tabular}
\end{tablebox}

\cmd{\%} は、ある数が偶数かどうか（2 で割ったあまりが 0 かどうか）を調べるときなどによく使います。

\subsection{コメント}

プログラムの中で \cmd{\#} を書くと、そこから行の終わりまでが\textbf{コメント}になり、実行されません。
コメントは、プログラムの説明を書いたり、一時的にある行を実行しないようにしたり（\textbf{コメントアウト}）するのに使います。

\begin{lstlisting}[style=python, caption={コメントの例}, label={lst:python-basics-comment}]
# 2 の 10 乗を計算する
print(2 ** 10)   # 行の途中からコメントを書くこともできる
# print(2 ** 12) この行はコメントアウトされているので実行されない
\end{lstlisting}

\cmd{print( )} は、カッコの中の値を画面に表示する\textbf{関数}です（\ref{sec:python-basics-function}節）。
スクリプトファイルで実行するときは、\cmd{print} を使わないと計算結果は表示されません。

\begin{exercise}
\begin{enumerate}
  \item テンパズルに挑戦しましょう。\cmd{5}, \cmd{8}, \cmd{6}, \cmd{2} の 4 つの数字をこの順に並べ、間に演算子（\cmd{+}, \cmd{-}, \cmd{*}, \cmd{/}）とカッコを入れて、計算結果が 10 になる式を作ってください。
  \item ロボットが 200 秒間走りました。200 秒は何分何秒でしょうか。\cmd{//} と \cmd{\%} を使って計算してください。
\end{enumerate}
\end{exercise}

\section{変数とデータ型}
\label{sec:python-basics-variable}

\subsection{変数}

\textbf{変数}は、値に名前を付けて保存しておくしくみです。
名前の付いた箱に値を入れておくイメージです。
\cmd{変数名 = 値} と書くと、変数に値が入ります（\textbf{代入}）。

\begin{lstlisting}[style=python, caption={変数を使う（speed.py）}, label={lst:python-basics-var}]
distance = 12.0   # 走った距離 [m]
time = 30.0       # かかった時間 [s]
speed = distance / time
print(speed)
\end{lstlisting}

\begin{terminal}[変数を使うプログラムを実行する]
$ python3 speed.py
0.4
\end{terminal}

変数の名前は、何の値が入っているのかがわかるように付けましょう。
Python では、単語をアンダースコアでつなぐ \cmd{max\_speed} のような名前が一般的です（\ref{sec:python-class-style}節）。

プログラムの \cmd{=} は、数学の「等しい」ではなく、「右辺の値を左辺の変数に入れる」という意味です。
そのため、変数の値を使って計算し、同じ変数に入れ直すこともできます。
\cmd{+=} などの演算子を使うと、これを短く書けます。

\begin{lstlisting}[style=python, caption={変数の値を更新する}, label={lst:python-basics-update}]
count = 0
count = count + 1   # count の値に 1 を足して、count に入れ直す
count += 1          # count = count + 1 と同じ意味
print(count)        # 2 と表示される
\end{lstlisting}

\subsection{データ型}

値には、\textbf{データ型}（型）があります。
よく使う型は次のとおりです。
値の型は、\cmd{type( )} 関数で確かめられます。

\begin{tablebox}[よく使うデータ型][tab:python-basics-types]
  \begin{tabular}{llp{0.45\linewidth}}
    \toprule
    型 & 意味 & 例 \\
    \midrule
    \cmd{int} & 整数 & \cmd{10}, \cmd{-3} \\
    \cmd{float} & 小数（浮動小数点数） & \cmd{0.5}, \cmd{3.14} \\
    \cmd{str} & 文字列 & \cmd{"robot"}, \cmd{'Hello'} \\
    \cmd{bool} & 真偽値（\ref{sec:python-basics-if}節） & \cmd{True}, \cmd{False} \\
    \bottomrule
  \end{tabular}
\end{tablebox}

\begin{terminal}[値の型を調べる]
>>> type(10)
<class 'int'>
>>> type(0.5)
<class 'float'>
>>> type("robot")
<class 'str'>
\end{terminal}

このほかに、「値がない」ことを表す \cmd{None} という特別な値もあります（\ref{sec:python-basics-none}節）。

\subsection{文字列}

文字列は、\cmd{"}（ダブルクォート）か \cmd{'}（シングルクォート）で囲んで書きます。
文字列どうしは \cmd{+} でつなげられます。

\begin{terminal}[文字列をつなげる]
>>> name = "turtle"
>>> "Hello, " + name + "!"
'Hello, turtle!'
\end{terminal}

文字列の中に \cmd{'} を含めたいときは、\cmd{"} で囲むか、\cmd{\textbackslash '} のように \cmd{\textbackslash}（バックスラッシュ）を付けます。
このような書き方を\textbf{エスケープシーケンス}といいます。
\cmd{\textbackslash n} は改行を表します。

\begin{terminal}[エスケープシーケンス]
>>> print("It's my robot.")
It's my robot.
>>> print('It\'s my robot.')
It's my robot.
>>> print("1 行目\n2 行目")
1 行目
2 行目
\end{terminal}

\subsection{型の変換}

数値と文字列は、そのままでは \cmd{+} でつなげられず、エラーになります。

\begin{terminal}[数値と文字列はつなげられない]
>>> temperature = 20
>>> temperature + "℃"
Traceback (most recent call last):
  File "<stdin>", line 1, in <module>
TypeError: unsupported operand type(s) for +: 'int' and 'str'
\end{terminal}

型を変換するには、\cmd{str( )}（文字列にする）、\cmd{int( )}（整数にする）、\cmd{float( )}（小数にする）を使います。

\begin{terminal}[型を変換する]
>>> str(temperature) + "℃"
'20℃'
>>> int("42") + 1
43
>>> float("3.14")
3.14
\end{terminal}

\cmd{int("5.5")} や \cmd{int("にじゅう")} のように、変換できない文字列を渡すとエラーになります。

\subsection{f 文字列}

変数の値を文字列の中に埋め込むときは、\textbf{f 文字列}を使うと便利です。
文字列の前に \cmd{f} を付け、埋め込みたい変数や式を \cmd{\{ \}} で囲みます。
数値も自動で文字列に変換されるので、\cmd{str( )} を使う必要がありません。
\cmd{:.2f} のように書くと、小数点以下の桁数を指定できます。

\begin{lstlisting}[style=python, caption={f 文字列を使う}, label={lst:python-basics-fstring}]
name = "turtle"
speed = 0.4
print(f"{name} の速度は {speed} m/s です")
print(f"{name} の速度は {speed * 3.6:.2f} km/h です")
\end{lstlisting}

\begin{terminal}[f 文字列の実行結果]
turtle の速度は 0.4 m/s です
turtle の速度は 1.44 km/h です
\end{terminal}

ROS 2 のプログラムでログを表示するときにも、f 文字列をよく使います（\ref{chap:rclpy}章）。

\subsection{キーボードからの入力}

\cmd{input( )} 関数を使うと、キーボードから入力された文字列を受け取れます。
カッコの中には、入力を促すメッセージを書きます。

\begin{lstlisting}[style=python, caption={入力を受け取る（greet.py）}, label={lst:python-basics-input}]
name = input("名前を入力してください=> ")
print(f"こんにちは、{name}さん!")
\end{lstlisting}

\begin{terminal}[入力を受け取るプログラムを実行する]
$ python3 greet.py
名前を入力してください=> ロボ太
こんにちは、ロボ太さん!
\end{terminal}

\cmd{input( )} で受け取った値は、数字を入力しても\textbf{文字列}になることに注意しましょう。
計算に使うときは、\cmd{int( )} や \cmd{float( )} で数値に変換します。

\begin{lstlisting}[style=python, caption={入力を数値に変換する}, label={lst:python-basics-input-num}]
radius = float(input("半径 [m] => "))
print(f"面積は {radius * radius * 3.14} m^2 です")
\end{lstlisting}

\begin{column}{小数の計算の誤差}
  コンピュータは、小数を 2 進数で近似して扱っています。
  そのため、\cmd{0.1 + 0.2} を計算すると \cmd{0.30000000000000004} になるように、小数の計算にはごくわずかな誤差が出ることがあります。
  表示するときは、\cmd{round(値, 桁数)} で丸めたり、f 文字列の \cmd{:.2f} で桁数を指定したりしましょう。
  また、小数どうしを \cmd{==}（\ref{sec:python-basics-if}節）で比べると、誤差のせいで等しくならないことがあるので注意が必要です。
\end{column}

\begin{exercise}
\begin{enumerate}
  \item 所持金が 4000 円あります。本（980 円、税率 10\%）を 2 冊と、ポテトチップス（250 円、税率 8\%）とアイス（150 円、税率 8\%）を 1 つずつ買った後の所持金を計算してください。それぞれの値を、何の値かがわかる名前の変数に入れてから計算すること。結果は \cmd{round( )} で整数に丸めて表示してください。
  \item \cmd{input( )} で見たことのあるアニメの題名を入力し、「〇〇は△文字です。」と表示してください。文字数は \cmd{len(文字列)} で調べられます。
  \item ロボットの車輪の半径 [m] と、1 秒あたりの車輪の回転数を \cmd{input( )} で入力し、ロボットが進む速さ [m/s]（$2 \times 3.14 \times \text{半径} \times \text{回転数}$）を、小数点以下 2 桁まで表示してください。
\end{enumerate}
\end{exercise}

\section{リスト・タプル・辞書}
\label{sec:python-basics-list}

\subsection{リスト}

\textbf{リスト}を使うと、複数の値を 1 つの変数にまとめて入れられます。
値をカンマで区切って \cmd{[ ]} で囲みます。
リストの中の 1 つ 1 つの値を\textbf{要素}といいます。
たとえば、センサで測った 1 週間分の距離のデータを、7 個の変数ではなく 1 つのリストで扱えます。

\begin{lstlisting}[style=python, caption={リストを作る}, label={lst:python-basics-list}]
distances = [1.2, 0.8, 0.5, 0.9, 1.5]   # センサで測った距離 [m]
print(distances)
print(len(distances))   # 要素の数
\end{lstlisting}

\begin{terminal}[リストの実行結果]
[1.2, 0.8, 0.5, 0.9, 1.5]
5
\end{terminal}

\subsection{要素を取り出す}

リストの要素には、先頭から順に\textbf{インデックス}（位置番号）が付いています。
インデックスは \textbf{1 からではなく 0 から}始まることに注意しましょう。
\cmd{リスト[インデックス]} で要素を取り出せます。
マイナスの数を指定すると、後ろから数えます（一番後ろは \cmd{-1}）。

\begin{terminal}[要素を取り出す]
>>> words = ["あ", "い", "う", "え", "お"]
>>> words[0]
'あ'
>>> words[3]
'え'
>>> words[-1]
'お'
\end{terminal}

\cmd{リスト[開始:終了]} と書くと、複数の要素をまとめて取り出せます（\textbf{スライス}）。
取り出されるのは、開始の位置から、終了の位置の\textbf{1 つ手前}までです。
開始を省略すると先頭から、終了を省略すると最後までになります。

\begin{terminal}[スライスで複数の要素を取り出す]
>>> words[1:3]
['い', 'う']
>>> words[:2]
['あ', 'い']
>>> words[2:]
['う', 'え', 'お']
\end{terminal}

\subsection{要素を変更・追加・削除する}

リストの要素は、後から変更したり、追加・削除したりできます。

\begin{tablebox}[リストの要素を変更・追加・削除する方法][tab:python-basics-list-methods]
  \begin{tabular}{lp{0.55\linewidth}}
    \toprule
    書き方 & 動作 \\
    \midrule
    \cmd{リスト[i] = 値} & i 番目の要素を変更する \\
    \cmd{リスト.append(値)} & 末尾に要素を追加する \\
    \cmd{リスト.insert(i, 値)} & i 番目に要素を挿入する \\
    \cmd{リスト.pop(i)} & i 番目の要素を取り出して削除する \\
    \cmd{リスト.remove(値)} & 指定した値の要素を削除する \\
    \bottomrule
  \end{tabular}
\end{tablebox}

\begin{terminal}[要素を変更・追加・削除する]
>>> sensors = ["camera", "lidar"]
>>> sensors.append("imu")
>>> sensors.insert(0, "gps")
>>> sensors
['gps', 'camera', 'lidar', 'imu']
>>> sensors.remove("camera")
>>> sensors
['gps', 'lidar', 'imu']
\end{terminal}

\cmd{append} のように、\cmd{値.名前( )} の形で使う関数を\textbf{メソッド}といいます（\ref{sec:python-class-class}節）。
数値のリストでは、\cmd{sum( )}（合計）、\cmd{max( )}（最大値）、\cmd{min( )}（最小値）も便利です。

\begin{terminal}[数値のリストの合計・最大値・最小値]
>>> sum(distances)
4.9
>>> min(distances)
0.5
\end{terminal}

\subsection{タプル}

\textbf{タプル}は、リストと似ていますが、作った後で要素を変更できないデータです。
値をカンマで区切って \cmd{( )} で囲みます。
ロボットの位置 \cmd{(x, y)} のように、組にして扱う値によく使います。
タプルの要素は、複数の変数に一度に取り出せます。

\begin{terminal}[タプルを使う]
>>> position = (1.5, -0.3)
>>> position[0]
1.5
>>> x, y = position
>>> y
-0.3
\end{terminal}

\subsection{辞書}

\textbf{辞書}は、\textbf{キー}と\textbf{値}の組を集めたデータです。
インデックスの代わりに、キー（名前）で値を取り出せます。
\cmd{\{キー: 値, キー: 値\}} の形で書きます。
ロボットの設定値（パラメータ）のように、名前で値を管理したいときに便利です。

\begin{lstlisting}[style=python, caption={辞書を使う}, label={lst:python-basics-dict}]
params = {"name": "turtle", "max_speed": 0.5, "wheel_radius": 0.033}
print(params["max_speed"])     # キーで値を取り出す
params["max_speed"] = 0.3      # 値を変更する
params["use_lidar"] = True     # 新しいキーと値を追加する
print(params)
\end{lstlisting}

\begin{terminal}[辞書の実行結果]
0.5
{'name': 'turtle', 'max_speed': 0.3, 'wheel_radius': 0.033, 'use_lidar': True}
\end{terminal}

辞書にないキーを指定するとエラーになります。
\cmd{"max\_speed" in params} と書くと、そのキーがあるかどうかを確かめられます。

\begin{exercise}
\begin{enumerate}
  \item \cmd{"green"}, \cmd{"blue"}, \cmd{"red"}, \cmd{"black"}, \cmd{"white"} の 5 つの要素を持つリストを作り、スライスを使って \cmd{['blue', 'red', 'black']} と表示してください。
  \item \cmd{["satou", "tanaka", "suzuki"]} というリストに \cmd{"itou"} を追加し、\cmd{"satou"} を削除してから、リスト全体とインデックス 1 の要素を表示してください。
  \item 3 人の身長 [cm] のリスト \cmd{[176, 158, 164]} から、平均の身長を計算して表示してください。
  \item ロボットの設定を表す辞書 \cmd{\{"name": "turtle", "max\_speed": 0.5\}} を作り、\cmd{max\_speed} を 0.8 に変更し、\cmd{"battery": 12.0} を追加して、「turtle の最高速度は 0.8 m/s」と表示してください。
\end{enumerate}
\end{exercise}

\section{条件分岐}
\label{sec:python-basics-if}

\subsection{if 文}

\textbf{if 文}を使うと、条件によって実行する処理を変えられます。

\begin{lstlisting}[style=python, caption={if 文の書き方}, label={lst:python-basics-if-syntax}]
if 条件A:
    条件A が成り立つときの処理
elif 条件B:
    条件A が成り立たず、条件B が成り立つときの処理
else:
    どの条件も成り立たないときの処理
\end{lstlisting}

条件の後ろには \cmd{:}（コロン）を付け、その下の処理は\textbf{インデント}（行頭の字下げ）をして書きます。
Python では、インデントによって「どこからどこまでが if 文の中の処理か」を表します（\ref{sec:python-features}節）。
インデントは半角スペース 4 つが標準です。
\cmd{elif} はいくつでも書くことができ、\cmd{elif} と \cmd{else} は必要なければ省略できます。

\begin{lstlisting}[style=python, caption={障害物までの距離で動作を変える（obstacle.py）}, label={lst:python-basics-if}]
distance = float(input("障害物までの距離 [m] => "))

if distance < 0.3:
    print("停止します")
elif distance < 1.0:
    print("減速します")
else:
    print("そのまま進みます")
\end{lstlisting}

\begin{terminal}[if 文のプログラムを実行する]
$ python3 obstacle.py
障害物までの距離 [m] => 0.5
減速します
\end{terminal}

\subsection{比較演算子}

条件は、主に次の\textbf{比較演算子}で書きます。
比較の結果は、成り立つときは \cmd{True}（真）、成り立たないときは \cmd{False}（偽）になります。

\begin{tablebox}[比較演算子][tab:python-basics-comparison]
  \begin{tabular}{ll}
    \toprule
    演算子 & 意味 \\
    \midrule
    \cmd{a == b} & a と b が等しい \\
    \cmd{a != b} & a と b が等しくない \\
    \cmd{a < b} / \cmd{a > b} & a が b より小さい / 大きい \\
    \cmd{a <= b} / \cmd{a >= b} & a が b 以下 / 以上 \\
    \bottomrule
  \end{tabular}
\end{tablebox}

「等しい」は \cmd{=} ではなく \cmd{==} と書くことに注意しましょう（\cmd{=} は代入です）。

\subsection{論理演算子}

複数の条件を組み合わせるときは、\textbf{論理演算子}を使います。

\begin{tablebox}[論理演算子][tab:python-basics-logical]
  \begin{tabular}{ll}
    \toprule
    演算子 & 意味 \\
    \midrule
    \cmd{a and b} & a と b がどちらも成り立つ \\
    \cmd{a or b} & a と b の少なくとも一方が成り立つ \\
    \cmd{not a} & a が成り立たない \\
    \bottomrule
  \end{tabular}
\end{tablebox}

\begin{terminal}[論理演算子]
>>> battery = 11.5
>>> 11.0 <= battery and battery <= 12.6
True
>>> battery < 11.0 or battery > 12.6
False
\end{terminal}

\begin{caution}[\cmd{or} の書き間違い]
  「\cmd{x} を 5 で割ったあまりが 2 か 3」という条件を、\cmd{x \% 5 == 2 or 3} と書いてはいけません。
  これは \cmd{(x \% 5 == 2) or 3} と解釈され、\cmd{3} は常に真とみなされるので、条件がいつも成り立ってしまいます。
  正しくは \cmd{x \% 5 == 2 or x \% 5 == 3} と書きます。
\end{caution}

\subsection{値がないことを表す None}
\label{sec:python-basics-none}

プログラムでは、「まだ値が決まっていない」「結果がない」ことを表したい場面がよくあります。
たとえば、センサのデータを受け取るプログラムでは、最初のデータが届くまでは、使える値がありません。
このようなときは、\textbf{None} という特別な値を使います。
\cmd{None} は \cmd{0} や空の文字列 \cmd{""} とは違い、「何もない」ことを表す値で、型は \cmd{NoneType} です。

変数が \cmd{None} かどうかを調べるときは、\cmd{==} ではなく \cmd{is} を使って、\cmd{x is None}（\cmd{None} である）、\cmd{x is not None}（\cmd{None} でない）と書きます。

\begin{lstlisting}[style=python, caption={None で「まだ値がない」ことを表す}, label={lst:python-basics-none}]
distance = None   # まだ距離を測っていない

if distance is None:
    print("まだ距離のデータがありません")
else:
    print(f"距離は {distance} m です")
\end{lstlisting}

\begin{terminal}[None を使ったプログラムの実行結果]
まだ距離のデータがありません
\end{terminal}

1 行目を \cmd{distance = 0.8} に変えて実行すると、「距離は 0.8 m です」と表示されます。
ここで \cmd{distance = 0} として「まだ測っていない」を表すと、本当に距離が 0 m だったときと区別できなくなります。
\cmd{None} を使えば、その心配がありません。

\cmd{None} は、\ref{sec:python-basics-function}節で学ぶ関数で、結果を返せないときの戻り値としてもよく使われます。
\ref{chap:rclpy}章で ROS 2 のノードを書くときも、「まだメッセージが届いていない」ことを \cmd{None} で表します。

\begin{exercise}
\begin{enumerate}
  \item \cmd{input( )} で 2 つの数値 a, b を入力し、「a の方が大きい」「b の方が大きい」「a と b は同じ」のいずれかを表示してください。
  \item 身長 [m] と体重 [kg] を入力して BMI（体重 ÷ 身長の 2 乗）を計算し、18.5 未満なら「やせ」、25 未満なら「標準」、30 未満なら「肥満」、それ以上なら「高度肥満」と表示してください。
  \item 入力した整数が偶数なら「偶数である」、さらにそれが 25 以上なら「25 以上である」と表示してください。また、偶数かどうかに関係なく、3 の倍数なら「3 の倍数である」と表示してください。
  \item ロボットのバッテリの電圧を入力し、12.0 V 以上なら「充電十分」、11.0 V 以上 12.0 V 未満なら「そろそろ充電」、11.0 V 未満なら「すぐに充電してください」と表示してください。ただし、0 V 以下や 15 V より大きい値が入力されたときは「センサの値が異常です」と表示すること。
\end{enumerate}
\end{exercise}

\section{繰り返し}
\label{sec:python-basics-loop}

\subsection{for 文}

同じ処理を繰り返すには、\textbf{for 文}を使います。
\cmd{for 変数 in リストなど:} と書くと、リストなどの要素が先頭から 1 つずつ変数に入り、その下のインデントした処理が繰り返されます。

\begin{lstlisting}[style=python, caption={for 文でリストの要素を順に処理する}, label={lst:python-basics-for}]
sensors = ["camera", "lidar", "imu"]
for sensor in sensors:
    print(f"{sensor} を確認します")
\end{lstlisting}

\begin{terminal}[for 文の実行結果]
camera を確認します
lidar を確認します
imu を確認します
\end{terminal}

決まった回数だけ繰り返すときは、\cmd{range( )} を使います。
\cmd{range(5)} は 0 から 4 までの 5 つの整数を、\cmd{range(1, 6)} は 1 から 5 までの整数を順に作ります（終わりの数は含まれません）。

\begin{lstlisting}[style=python, caption={range で決まった回数だけ繰り返す}, label={lst:python-basics-range}]
total = 0
for i in range(1, 11):
    total += i
print(total)   # 1 から 10 までの合計 55 が表示される
\end{lstlisting}

\subsection{break と continue}

\cmd{break} を実行すると、その時点で繰り返しを終了します。
\cmd{continue} を実行すると、それより後の処理を飛ばして、次の繰り返しに進みます。

\begin{lstlisting}[style=python, caption={break と continue}, label={lst:python-basics-break}]
distances = [1.2, -1.0, 0.8, 0.2, 1.5]   # -1.0 は測定に失敗した値
for d in distances:
    if d < 0:
        continue          # 失敗した値は飛ばす
    if d < 0.3:
        print(f"{d} m: 障害物が近いので停止")
        break             # 繰り返しを終了する
    print(f"{d} m: 前進")
\end{lstlisting}

\begin{terminal}[break と continue の実行結果]
1.2 m: 前進
0.8 m: 前進
0.2 m: 障害物が近いので停止
\end{terminal}

\subsection{while 文}

\textbf{while 文}は、条件が成り立っている間、処理を繰り返します。
何回繰り返すかが前もって決まっていないときに便利です。

\begin{lstlisting}[style=python, caption={while 文で繰り返す}, label={lst:python-basics-while}]
battery = 12.6
minutes = 0
while battery > 11.0:
    battery -= 0.2      # 1 分ごとに 0.2 V 下がるとする
    minutes += 1
print(f"{minutes} 分で充電が必要になります")
\end{lstlisting}

\cmd{while True:} と書くと、\cmd{break} するまで永遠に繰り返します。
入力が正しい値になるまで聞き直す、といった処理によく使います。

\begin{lstlisting}[style=python, caption={正しい値が入力されるまで聞き直す}, label={lst:python-basics-while-true}]
while True:
    num = int(input("1〜100 の値を入力してください=> "))
    if 1 <= num <= 100:
        break
    print("範囲外の値です")
print(f"入力された値は {num} です")
\end{lstlisting}

\cmd{1 <= num <= 100} のように、比較演算子をつなげて範囲を表すこともできます。

\begin{caution}[無限ループ]
  条件を間違えると、繰り返しが永遠に終わらない\textbf{無限ループ}になってしまいます。
  プログラムが止まらなくなったら、\key{Ctrl}+\key{C} で強制終了しましょう（\ref{sec:process-kill}節）。
\end{caution}

\subsection{リスト内包表記}
\label{sec:python-basics-comprehension}

リストの要素を 1 つずつ処理して、新しいリストを作りたいことがあります。
for 文で書くと、次のようになります。

\begin{lstlisting}[style=python, caption={for 文で新しいリストを作る}, label={lst:python-basics-for-list}]
distances = [1.2, -1.0, 0.8, 0.2, 1.5]   # [m]、-1.0 は測定に失敗した値
distances_cm = []
for d in distances:
    distances_cm.append(d * 100)
print(distances_cm)   # [120.0, -100.0, 80.0, 20.0, 150.0]
\end{lstlisting}

Python には、これを 1 行で書ける\textbf{リスト内包表記}という書き方があります。
\cmd{[式 for 変数 in リスト]} と書くと、リストの要素を 1 つずつ変数に入れて式を計算し、その結果を並べた新しいリストができます。
後ろに \cmd{if 条件} を付けると、条件が成り立つ要素だけを取り出せます。

\begin{lstlisting}[style=python, caption={リスト内包表記}, label={lst:python-basics-comprehension}]
distances = [1.2, -1.0, 0.8, 0.2, 1.5]

distances_cm = [d * 100 for d in distances]
print(distances_cm)   # [120.0, -100.0, 80.0, 20.0, 150.0]

valid = [d for d in distances if d >= 0]   # 正しく測れた値だけ
print(valid)          # [1.2, 0.8, 0.2, 1.5]
print(min(valid))     # 一番近い障害物までの距離 0.2
\end{lstlisting}

LiDAR（レーザで周りの物までの距離を測るセンサ）の距離のデータのように、たくさんの値の中から必要なものだけを取り出す処理は、ロボットのプログラムでよく出てきます。
最初は for 文で書いても構いませんが、ほかの人のコードを読むときのために、この書き方も覚えておきましょう。

\subsection{乱数を使う}

練習問題でゲームを作るために、\textbf{乱数}（でたらめな数）の作り方を紹介しておきます。
\cmd{import random} と書いておくと、\cmd{random.randint(a, b)} で a 以上 b 以下の整数の乱数を作れます。
\cmd{import} については、\ref{sec:python-class-module}節で詳しく学びます。

\begin{lstlisting}[style=python, caption={サイコロを 5 回振る}, label={lst:python-basics-random}]
import random

for i in range(5):
    print(random.randint(1, 6))
\end{lstlisting}

\begin{exercise}
\begin{enumerate}
  \item for 文を使って、1 から 30 までの整数の合計を計算して表示してください。
  \item リスト \cmd{[1, 62, 3, 6743, 51, 6783, 42, 4671, 5, 236]} から、1000 未満の値だけを表示してください。\cmd{continue} を使い、\cmd{print} は if 文の中に書かないこと。
  \item 整数を入力して、1 からその数までの棒グラフを表示してください（入力が 4 なら、下のように表示します）。\cmd{print("■", end="")} と書くと、改行せずに表示できます。また、\cmd{"■" * 3} のように文字列に整数を掛けると、文字列を繰り返せます。
\begin{verbatim}
1:■
2:■■
3:■■■
4:■■■■
\end{verbatim}
  \item 数を次々と入力し、負の数が入力されたら、それまでに入力された数の個数・合計・平均を表示してください。
  \item 数当てゲームを作りましょう。1 から 30 までの乱数を作り、5 回以内に当ててもらいます。外れたときは「それよりも大きいです」「それよりも小さいです」とヒントを出し、当たったら「正解です！」と表示して終了します。5 回で当てられなかったときは、正解を表示してください。
  \item じゃんけんのプログラムを作りましょう。グーを 0、チョキを 1、パーを 2 として自分の手を入力し、コンピュータの手を乱数で決めて、勝ち・負け・あいこを表示してください。ヒント：\cmd{(自分の手 - 相手の手 + 3) \% 3} を計算すると、あいこなら 0、負けなら 1、勝ちなら 2 になります。
\end{enumerate}
\end{exercise}

\section{関数}
\label{sec:python-basics-function}

\subsection{関数とは}

\textbf{関数}は、ひとまとまりの処理に名前を付けたものです。
これまでに使ってきた \cmd{print( )}, \cmd{len( )}, \cmd{input( )} なども、Python に最初から用意されている関数（\textbf{組み込み関数}）です。
自分で関数を作っておけば、同じ処理を何度も書かずに、名前を呼ぶだけで実行できます。

\subsection{関数を作る}

関数は、\cmd{def} を使って作ります（\textbf{定義}します）。

\begin{lstlisting}[style=python, caption={関数の書き方}, label={lst:python-basics-def-syntax}]
def 関数名(引数1, 引数2, ...):
    処理
    return 戻り値
\end{lstlisting}

\textbf{引数}（ひきすう）は、関数に渡す値で、関数の中では変数として使えます。
\textbf{戻り値}は、関数が処理した結果として呼び出し元に返す値で、\cmd{return} で返します。
引数も戻り値も、必要なければ省略できます。
関数は、定義した後で \cmd{関数名(引数)} と書いて呼び出します。

\begin{lstlisting}[style=python, caption={円の面積を計算する関数（circle.py）}, label={lst:python-basics-def}]
def circle_area(radius):
    area = radius * radius * 3.14
    return area


a = circle_area(2.0)
print(a)
print(circle_area(5.0))
\end{lstlisting}

\begin{terminal}[関数のプログラムを実行する]
$ python3 circle.py
12.56
78.5
\end{terminal}

関数の中の処理は、呼び出されるまで実行されません。
また、関数は呼び出すより前（ファイルの上のほう）で定義しておく必要があります。

\subsection{いろいろな引数と戻り値}

引数には、\cmd{引数名=値} の形で\textbf{デフォルト値}（省略したときに使われる値）を決めておけます。
呼び出すときに \cmd{引数名=値} と書くと（\textbf{キーワード引数}）、引数の順番を気にせずに値を渡せます。
また、\cmd{return} の後ろにカンマで区切って書くと、複数の値を返せます（実際にはタプルとして返されます）。

\begin{lstlisting}[style=python, caption={デフォルト値と複数の戻り値}, label={lst:python-basics-def2}]
def price_with_tax(price, tax_rate=0.1):
    return round(price * (1 + tax_rate))


def min_max(values):
    return min(values), max(values)


print(price_with_tax(980))                  # 税率 10% で計算
print(price_with_tax(250, tax_rate=0.08))   # 税率 8% で計算
low, high = min_max([1.2, 0.8, 0.5, 1.5])
print(low, high)
\end{lstlisting}

\begin{terminal}[デフォルト値と複数の戻り値の実行結果]
1078
270
0.5 1.5
\end{terminal}

\cmd{print( )} にカンマで区切って複数の値を渡すと、空白で区切って表示されます。

\subsection{変数の有効範囲}

関数の中で作った変数や引数は、その関数の中でしか使えません。
関数の外でその変数を使おうとすると、\cmd{NameError}（その名前は定義されていない）というエラーになります。
関数の中で計算した結果を外で使いたいときは、\cmd{return} で返します。

\begin{terminal}[関数の中の変数は外から使えない]
>>> def f(x):
...     temp = x * 2
...     return temp
...
>>> f(3)
6
>>> temp
Traceback (most recent call last):
  File "<stdin>", line 1, in <module>
NameError: name 'temp' is not defined
\end{terminal}

\begin{exercise}
\begin{enumerate}
  \item 「完璧を」「目指すよりも」「まずは」「終わらせろ」の 4 行を表示する関数を作り、3 回呼び出してください（ループは使わないこと）。
  \item 整数を 1 つ受け取り、偶数なら「偶数」、奇数なら「奇数」と表示する関数を作り、3 回以上呼び出してください。
  \item 3 つの数値を受け取り、その中で最も大きい数値を返す関数を作ってください。\cmd{max( )} を使わないこと。表示は関数の外で行うこと。
  \item リスト \cmd{[3, 62, 4, 56, 346, 8, 2, 893, 52, 46]} を受け取り、最大値と最小値の 2 つを返す関数を作ってください。表示は関数の外で行うこと。
  \item 距離のリストを受け取り、0 以上の値（正しく測れた値）だけの平均を返す関数を作ってください。正しい値が 1 つもないときは \cmd{None}（値がないことを表す特別な値）を返すこと。
\end{enumerate}
\end{exercise}

\section*{この章のまとめ}

\begin{itemize}
  \item 数値は \cmd{+}, \cmd{-}, \cmd{*}, \cmd{/}, \cmd{**}, \cmd{//}, \cmd{\%} で計算できます。\cmd{\#} から後ろはコメントです。
  \item 変数には \cmd{=} で値を代入します。値には \cmd{int}, \cmd{float}, \cmd{str}, \cmd{bool} などの型があり、\cmd{int( )}, \cmd{float( )}, \cmd{str( )} で変換します。変数の値を文字列に埋め込むには f 文字列が便利です。
  \item \cmd{input( )} で受け取った値は文字列なので、計算に使うときは数値に変換します。
  \item リストは \cmd{[ ]}、タプルは \cmd{( )}、辞書は \cmd{\{ \}} で作ります。リストのインデックスは 0 から始まります。
  \item if 文で条件分岐、for 文と while 文で繰り返しができます。処理のまとまりはインデントで表します。
  \item \cmd{def} で関数を定義し、引数で値を受け取り、\cmd{return} で結果を返します。
\end{itemize}
